% ============================================================================
% Fichier  : partie1/P01_C03_gouvernance_indicateurs.tex
% Titre    : Gouvernance et Indicateurs de la Supervision
% Auteur   : MINKA MI NGUIDJOÏ Thierry Emmanuel
% ============================================================================
% Sources : ISO/IEC 27004 (mesure) ; ISO/IEC 27035 ; CIS Controls v8.
% ============================================================================

\chapter{Gouvernance et Indicateurs de la Supervision}
\label{chap:gouvernance}
\chaptermark{Gouvernance et indicateurs}

\epigraph{%
    « Ce qui ne se mesure pas ne se pilote pas~;\\
    ce qui ne se pilote pas se dégrade. »%
}{Principe du cours}

\niveauChapitre{Fondamental}{Chapitres~\ref{chap:pourquoi} et~\ref{chap:soc}.}

\begin{boxobjectifs}
\begin{itemize}
    \item Définir la gouvernance d'un dispositif de supervision.
    \item Distinguer indicateurs d'activité, de performance et de résultat.
    \item Construire un tableau de bord utile et honnête.
    \item Fixer des objectifs de délai (détection, réaction).
    \item Rendre compte à la direction sans noyer sous la technique.
\end{itemize}
\end{boxobjectifs}

\bigskip

% ============================================================================
\section*{Question directrice : comment savoir si la supervision de LiZbiZ s'améliore vraiment~?}
% ============================================================================

Après l'incident, LiZbiZ met en place une détection. Six mois plus tard,
M.~ETEME demande~: « Sommes-nous mieux protégés~? » La réponse ne peut être
ni « oui » vague ni une avalanche de courbes~: elle tient dans quelques
indicateurs choisis et suivis dans le temps.

\medskip
\textbf{Plan du chapitre.}
\begin{enumerate}
    \item Gouverner la supervision (section~\ref{sec:c03-gouverner}).
    \item Trois familles d'indicateurs (section~\ref{sec:c03-familles}).
    \item Le tableau de bord (section~\ref{sec:c03-tdb}).
    \item Les objectifs de délai (section~\ref{sec:c03-delais}).
    \item Rendre compte (section~\ref{sec:c03-rendre}).
\end{enumerate}

\separateur

% ============================================================================
\section{Gouverner la Supervision}
\label{sec:c03-gouverner}
% ============================================================================

\begin{boxdef}
\termecle{Gouvernance de la supervision}~: ensemble des décisions qui fixent
le périmètre surveillé, les priorités, les moyens, les responsabilités et les
indicateurs, et qui en rendent compte à la direction. Elle relie la
technique à la stratégie.
\end{boxdef}

% ============================================================================
\section{Trois Familles d'Indicateurs}
\label{sec:c03-familles}
% ============================================================================

\begin{tableaucours}{L{3.2cm}L{4.6cm}L{4.4cm}}{%
    \tableauHeader{Famille} &
    \tableauHeader{Question} &
    \tableauHeader{Exemple}
}
Activité & Que fait-on~? & Nombre d'alertes traitées \\
Performance & Le fait-on bien et vite~? & Délai moyen de détection \\
Résultat & Est-on plus sûr~? & Part des techniques détectées, incidents évités \\
\end{tableaucours}
\vspace{-0.8em}
\begin{center}{\small\itshape Tableau~3.1. Trois familles d'indicateurs~\cite{iso27004}}\end{center}

\begin{boxalert}
\textbf{Erreur fréquente.} Ne suivre que l'activité (« nous avons traité
3~000 alertes »).

\smallskip
\textbf{Règle~:} l'activité ne dit rien de l'efficacité. Un volume d'alertes
élevé peut signaler de mauvaises règles, non une bonne détection. On
privilégie la performance et le résultat.
\end{boxalert}

% ============================================================================
\section{Le Tableau de Bord}
\label{sec:c03-tdb}
% ============================================================================

\begin{boxprocedure}
\textbf{Construire un tableau de bord de supervision}
\begin{enumerate}
    \item Choisir cinq à sept indicateurs, pas plus.
    \item Au moins un de chaque famille (activité, performance, résultat).
    \item Pour chacun~: définition, source, cible, tendance.
    \item Montrer l'évolution, pas seulement la valeur du mois.
    \item Dire honnêtement ce qui n'est pas couvert.
\end{enumerate}
\end{boxprocedure}

% ============================================================================
\section{Les Objectifs de Délai}
\label{sec:c03-delais}
% ============================================================================

Les deux délais cardinaux sont le \termecle{délai de détection} et le
\termecle{délai de réaction}. On les fixe par type d'incident~: un
rançongiciel exige une détection en minutes et un confinement en moins d'une
heure~; une anomalie mineure tolère un traitement en jours. Ces cibles
guident le dimensionnement du SOC (Chapitre~\ref{chap:soc}).

% ============================================================================
\section{Rendre Compte}
\label{sec:c03-rendre}
% ============================================================================

\begin{boxscenario}{Opération BÉNOUÉ. Six mois après}
\textbf{Constatation.} Nouveau tableau de bord~: délai de détection médian
ramené à 2~heures, couverture de journalisation portée à 85~\%, 40~\% des
techniques ATT\&CK jouées en test désormais détectées.

\medskip
\textbf{Analyse} (confiance moyenne). Les chiffres montrent un progrès réel,
mais la couverture reste partielle (15~\% des systèmes non journalisés) et la
détection en direct limitée~: ce sont les NEK à traiter
(Chapitre~\ref{chap:mesurer}). Rendre compte, c'est aussi dire cela.
\end{boxscenario}

\begin{boxmetier}{RSSI et analyste SOC}
\textbf{Un tableau de bord par lecteur.} Le RSSI présente à la direction trois
indicateurs de résultat~; l'analyste suit en interne le détail de la
performance. Les deux disent la même vérité, à deux niveaux.
\end{boxmetier}

\separateur

\begin{boxresume}
\textbf{Ce qu'il faut retenir de ce chapitre~:}
\begin{itemize}
    \item \textbf{Gouverner}, c'est relier la supervision à la stratégie.
    \item \textbf{Activité, performance, résultat}~: trois familles.
    \item \textbf{L'activité seule} trompe~; on vise performance et résultat.
    \item \textbf{Cinq à sept indicateurs}, avec tendance et angles morts.
    \item \textbf{Des objectifs de délai} par type d'incident.
\end{itemize}
\end{boxresume}

\begin{boxinfo}
\textbf{Points de contrôle de l'élève.} Je sais~: définir la gouvernance de
la supervision~; classer un indicateur~; bâtir un tableau de bord~; fixer des
délais cibles~; rendre compte à deux niveaux.
\end{boxinfo}

\bigskip

\begin{boxexo}[title={\ding{45}\quad Exercice~3.1. Familles \niveaubase}]
Classez~: nombre d'alertes~; délai de détection~; part des techniques
détectées.
\end{boxexo}

\begin{boxexo}[title={\ding{45}\quad Exercice~3.2. Activité trompeuse \niveaubase}]
Pourquoi « 3~000 alertes traitées » ne prouve-t-il pas une bonne détection~?
\end{boxexo}

\begin{boxexo}[title={\ding{45}\quad Exercice~3.3. Combien \niveaubase}]
Combien d'indicateurs un tableau de bord devrait-il comporter, et pourquoi
pas plus~?
\end{boxexo}

\begin{boxexo}[title={\ding{45}\quad Exercice~3.4. Cibles \niveaumoyen}]
Proposez des délais cibles de détection et de réaction pour un rançongiciel,
et justifiez.
\end{boxexo}

\begin{boxexo}[title={\ding{45}\quad Exercice~3.5. Lire BÉNOUÉ \niveaumoyen}]
Commentez les chiffres de l'opération BÉNOUÉ à six mois~: progrès et limites.
\end{boxexo}

\begin{boxexo}[title={\ding{45}\quad Exercice~3.6. Angles morts \niveaumoyen}]
Pourquoi faut-il afficher ce qui n'est pas couvert dans un tableau de bord~?
\end{boxexo}

\begin{boxexo}[title={\ding{45}\quad Exercice~3.7. Deux lecteurs \niveaumoyen}]
Quels indicateurs montrez-vous à la direction~? Lesquels gardez-vous en
interne~?
\end{boxexo}

\begin{boxexo}[title={\ding{45}\quad Exercice~3.8. Tendance \niveaumoyen}]
Pourquoi la tendance importe-t-elle plus que la valeur d'un mois isolé~?
\end{boxexo}

\begin{boxexo}[title={\ding{45}\quad Exercice~3.9. Tableau de bord \niveauexpert}]
Construisez le tableau de bord complet de supervision de LiZbiZ~: sept
indicateurs, définition, source, cible.
\end{boxexo}

\begin{boxexo}[title={\ding{45}\quad Exercice~3.10. Présentation au comité \niveauexpert}]
Rédigez la page que le RSSI présente au comité de direction pour répondre à
« Sommes-nous mieux protégés~? ».
\end{boxexo}

\bigskip

\begin{boxinfo}
\textbf{Pour aller plus loin.}
\begin{itemize}
    \item ISO/IEC 27004~\cite{iso27004}~; ISO/IEC 27035~\cite{iso27035}.
\end{itemize}
\end{boxinfo}

\begin{boxinfo}
\textbf{Transition.} La supervision commence par les traces~: la Partie~II
s'ouvre sur les journaux (Chapitre~\ref{chap:journaux}).
\end{boxinfo}
